\par\medskip
\section{Guide to the individual calculations}\label{sec:individual}
The identifiers below connect the physical argument to a derivation, a formal or numerical check, and the archived inputs. Each entry describes a specified model or preparation. The electronic guide links directly to the files for that entry.

\subsection{How a gravitational contact can distinguish correlations}
\begin{description}\small
\setlength{\itemsep}{.25em}\setlength{\parskip}{0pt}
\item[\textnormal{\texttt{RES-0093}}] A circle-momentum scattering channel distinguishes recording transfer, its sign, and the full transferred integer. Their return probabilities have different large-radius limits even for the same incoming state.
\item[\textnormal{\texttt{RES-0094}}] Breakup of a bound gravitational mode supplies a continuous Born distribution. A single circular material reference can leave two positions indistinguishable; a second reference can remove that ambiguity at an energy cost.
\item[\textnormal{\texttt{RES-0095}}] An additional relative-curvature field can be dark to direct ordinary matter excitation. A static exterior giving the desired force does not in general suppress conversion of ordinary waves at both material ends.
\item[\textnormal{\texttt{RES-0096}}] A variance-dependent relative source shifts an instability threshold. A direct curvature interaction also requires local exchange and Bianchi consistency; equality of marginal variances does not guarantee equal interference.
\item[\textnormal{\texttt{RES-0097}}] Equal exterior mass and trapped area can coexist with different interior tidal records. A comparison can be silent on two identical unknown inputs; a uniformly normalized vertex can lose its effect on a decompactifying circle.
\item[\textnormal{\texttt{RES-0098}}] Five spectral samples can distinguish a direct chain from a circle having the same leading dispersion. Matrix gauge charge and soft gravitational radiation constrain which neutral material states can preserve the exchange phase.
\item[\textnormal{\texttt{RES-0099}}] A finite direct chain and a propagating field matched at zero frequency can respond at different early times. The work includes receiver recoil, and a returning material charge can have a finite lifetime.
\item[\textnormal{\texttt{RES-0100}}] Two models can match a four-point response and still deliver different work to rotating quadrupoles. Even membranes receiving equal work can subsequently respond with opposite signs.
\item[\textnormal{\texttt{RES-0101}}] Three correlated tidal inputs reveal a maximum relative action error of three fifths for the best specified additive-mediator approximation. Resolving the sign of a rapid matrix encounter also requires a finite angular reference.
\item[\textnormal{\texttt{RES-0102}}] Two matrix membranes set a finite limit on resolving a contact. Equal spectra for continuous junctions can leave different clock excitations, while a continuum of outgoing modes can remove a discrete projection.
\item[\textnormal{\texttt{RES-0103}}] A closed mediator cycle or a direct contact can imprint a phase without transferring tensor-wave occupation. Interference, recoil and time resolution constrain how much which-path information can remain.
\item[\textnormal{\texttt{RES-0104}}] Three independent receipts consume a physical budget. A positive mediator has a self-response as well as a mutual response; it cannot mimic every direct interaction at fixed cost.
\item[\textnormal{\texttt{RES-0105}}] A perfect exchange-parity record at a direct junction requires suitable exchange of the material paths. A finite-energy clock bounds how closely two receipts can be discriminated.
\item[\textnormal{\texttt{RES-0106}}] A naturally decaying probe can distinguish interactions with identical static force by residual work or recoverable parity. Two lifetimes bound the strength in the specified cascade.
\item[\textnormal{\texttt{RES-0107}}] Re-interference bounds how much charge can be written before two membranes reunite. Full D2/M2 transfer statistics can reject a model of independent jumps that fits selected mean rates.
\item[\textnormal{\texttt{RES-0108}}] A wrapped membrane gives a two-time receipt with a direct phase that must respect refinement of the layers. Four-dimensional two-form intersections and three-form flux ends have different local charge accounting.
\item[\textnormal{\texttt{RES-0109}}] A charged membrane ending a flux region can correlate received winding and continuous recoil. The resulting resolution depends on which quanta the receiver absorbs, even at the same temperature.
\end{description}

\subsection{Infall, reversal and optional sources of gravitational stress}
\begin{description}\small
\setlength{\itemsep}{.25em}\setlength{\parskip}{0pt}
\item[\textnormal{\texttt{RES-0110}}] A real scalar resonance fixes selection probabilities for a prepared quadrupole. In the specified isotropic band mismatch above one sixth lets the emitted channel carry reference entanglement.
\item[\textnormal{\texttt{RES-0111}}] The survival of an internal material record changes the threshold for entanglement transfer in a fixed outgoing wave packet. Natural fluorescence fixes an example of an absolute material rate.
\item[\textnormal{\texttt{RES-0112}}] A source-labelled octupole can deposit a quantum state in the surface channel. Positivity of three material tags, accessible species, and recoil jointly constrain how much can escape.
\item[\textnormal{\texttt{RES-0113}}] A pair-number model links material receipt to an expansion term proportional to the inverse sixth power of the scale factor. High-redshift light can reject branches still admitted by a distance-only fit.
\item[\textnormal{\texttt{RES-0114}}] A locally absorbing observer and an exterior surface can implement distinct continuations. A classical null surface offers no later local path to infinity; pre-encoding, a timelike surface or extra coupling is needed.
\item[\textnormal{\texttt{RES-0115}}] The tidal signal from an orbiting body can be recorded without resolving every material degree of freedom. Thermal noise, receiver frequency and recoil can prevent transfer of an unknown state even when energy arrives.
\item[\textnormal{\texttt{RES-0116}}] Integrating the relative work over many layers yields a finite-duration routing problem. The proper time to the classical centre and the exterior time available to encode outgoing quanta impose separate constraints.
\item[\textnormal{\texttt{RES-0117}}] Ordinary emission and gravitational exchange can populate many alternatives before arrival. Exact closure additionally requires an inward channel incapable of carrying the received relative charge.
\item[\textnormal{\texttt{RES-0118}}] Prepared thermal-light or cold-front conditions can suppress the remote inward probability for a human mass and a visible grain. The occupation and common near-field receipt are essential assumptions.
\item[\textnormal{\texttt{RES-0119}}] Common and difference charges from one spectrum can sharpen a local history while choosing a distant route. Reversing ordinary dynamics leaves a return probability determined by the charge groups that were actually selected.
\item[\textnormal{\texttt{RES-0120}}] A grain-sized receiver must pay preparation work and finish its physical reflection before the local deadline. A sharp four-way label alone does not guarantee a timely readout.
\item[\textnormal{\texttt{RES-0121}}] The same signed interface used in an infall experiment has an energy balance and can carry distinct stress components. Two physical charges constrain the form of a flux sector.
\item[\textnormal{\texttt{RES-0122}}] Material distributions giving pressure, optical distance and radial forces were compared to galaxy and expansion data. Certain dark-sector identifications fail while the black-hole mechanism can still be assessed independently.
\item[\textnormal{\texttt{RES-0123}}] A local gravitational record for a glass grain and correlated boundary pulses can alter routing. The bound from recoverable elastic strain is too small for this source alone to account for the inferred dark mass.
\end{description}

\subsection{The exterior boundary, finite time and inclusive radiation}
\begin{description}\small
\setlength{\itemsep}{.25em}\setlength{\parskip}{0pt}
\item[\textnormal{\texttt{RES-0124}}] A local incoming flux can supply finite gravitational charge. An empty port remains empty in the equal-stress control; energy and whole-channel recovery impose additional conditions on a complete evaporation history.
\item[\textnormal{\texttt{RES-0125}}] A finite-mass receiver converts a tidal impulse into recoil. Knowing the total emitted momentum still leaves multiplicity spaces at the same quantum numbers that can carry polarization information.
\item[\textnormal{\texttt{RES-0126}}] A finite receiving interval restricts the accessible charge. Exterior light-cone data constrain terminal Bondi mass, while a record of angular momentum can bound fidelity even at the same radiated energy.
\item[\textnormal{\texttt{RES-0127}}] A signed material record can be produced by curvature before arrival. A circuit of finite causal length can cancel a waveform and limits how much charge a contracting material ring can occupy.
\item[\textnormal{\texttt{RES-0128}}] A transported charge can thermalize or relax before it is read. Three-port recovery distinguishes the accessible record from an unobserved soft output and must respect terminal momentum conservation.
\item[\textnormal{\texttt{RES-0129}}] A local clock can leave too little time to prepare an outgoing pair, although a source is energetic enough. Integrated linewidth, propagation distance and recoil constrain a possible terminal channel.
\item[\textnormal{\texttt{RES-0130}}] The charge budget bounds how much of the receiving area can be closed in finite time. In a winding-protected entrance the result also depends on the local orientation of the deposited current.
\item[\textnormal{\texttt{RES-0131}}] Advance delivery of charge can release a material membrane's tension. Exterior two-point correlations depend on the physical capture radius and recoil in addition to ordinary luminosity.
\item[\textnormal{\texttt{RES-0132}}] A recoiling surface has a definite four-dimensional causal future in specified examples. Capture from a finite pulse and late emission must both fit the available propagation interval.
\item[\textnormal{\texttt{RES-0133}}] Finite support of an early charge and a local exchange into fusion states give a test of where the quantum memory resides. Transverse surface stability and the outgoing one-quantum spectrum constrain the model.
\item[\textnormal{\texttt{RES-0134}}] A curvature-sensitive surface may acquire relative charge only after integrating its tidal impulse. Capture phase along a short closed geodesic and finite absorption time restrict the useful surface region.
\item[\textnormal{\texttt{RES-0135}}] A falling tide can generate a signed defect before contact. Routing releases energy; two candidates with equal late luminosity can disagree on correlations between the ends of a circular memory.
\item[\textnormal{\texttt{RES-0136}}] A spherical dust source and a recoiling surface give a causal charge-arrival test. Matching radial energy flux and transporting a four-dimensional field are separate checks.
\item[\textnormal{\texttt{RES-0137}}] A retarded curvature source in dust can generate an oriented eigenline over the sphere. A sign detectable at one end does not certify an occupied compensating line, and pressure from the incoming and outgoing fields matters.
\item[\textnormal{\texttt{RES-0138}}] A recoiling material record formed in finite fall time constrains possible terminal two-photon recovery. The cumulative angular emission probability and unknown payload must be compared in the same channel.
\item[\textnormal{\texttt{RES-0139}}] A neutral singlet and a finite physical orientation reference can prepare boundary charge. Reabsorption of emitted photons changes the recovery of magnetic substates.
\item[\textnormal{\texttt{RES-0140}}] The charging time of a stellar surface can be compared directly with infall communication and local orbital response. Returning radiation can carry energy while reducing access to the original quantum coherence.
\item[\textnormal{\texttt{RES-0141}}] The finite spherical link reproduces the signed count of inward modes. Matching Hawking intensities for every open species leaves different answers for a coherent superposition, already before a quantum endpoint.
\end{description}

\subsection{The matrix action, material clocks and the observable geometry}
\begin{description}\small
\setlength{\itemsep}{.25em}\setlength{\parskip}{0pt}
\item[\textnormal{\texttt{RES-0142}}] A material spectral crossing in matrix dynamics supplies a candidate compensated charge. An energy-matched local instrument can route the payload and reduce to the specified charged-interface constraints.
\item[\textnormal{\texttt{RES-0143}}] Support domination limits exact route separation on a common port. Matched charge projectors and recoil give sharper criteria than finite rates. Spherical two-derivative gravity imposes an additional area-action test.
\item[\textnormal{\texttt{RES-0144}}] The adopted ordinary parent is a single IKKT clock sector with BFSS time evolution. A band-limited circle aperture leaks on every open angular arc; orthogonal integer charges remain available.
\item[\textnormal{\texttt{RES-0145}}] Ten target components follow with the specified ordinary multiplet. Tidal deadlines, spin-dependent charge probabilities, flagged recovery and a finite material-clock formulation can be computed in the same framework.
\item[\textnormal{\texttt{RES-0146}}] A material crossing can emerge from initially localized matrix matter. Axial and oriented receivers discriminate different spin information; endpoint arrival restricts the actual area that can be closed.
\item[\textnormal{\texttt{RES-0147}}] Vacuum fluctuations and prepared collective motion offer different sources for a receiver. A nine-colour preparation changes its index before the deadline, although its prepared motion also drives a crossing in the flat control.
\item[\textnormal{\texttt{RES-0148}}] Closed-cycle tests separate different physical resolutions of the same charge distribution. A family of ordinary matrix solutions realizes exact four-dimensional vacuum waves, including both tidal polarizations.
\item[\textnormal{\texttt{RES-0149}}] A dust-driven relational clock admits a stated window of initial relative velocities for a signed crossing. Co-falling zero-relative-velocity preparations give counterexamples to a universal trigger.
\item[\textnormal{\texttt{RES-0150}}] Matrices can give directional expansion and a receiver-dependent conditional distribution of geometries. A local supernova fit and radial BAO contrast test specific low-energy branches without counting all vacua.
\item[\textnormal{\texttt{RES-0151}}] Two inequivalent curved classical configurations suffice to establish a classical landscape in the sense used here. This does not supply a count or a probability measure on quantum vacua.
\end{description}

\subsection{Curvature, communication between observers and material reception}
\begin{description}\small
\setlength{\itemsep}{.25em}\setlength{\parskip}{0pt}
\item[\textnormal{\texttt{RES-0152}}] A short interval of strong curvature leads to a matrix evolution problem with a finite clock sector. Equal incoming tidal information is insufficient to infer a unique future quantum state.
\item[\textnormal{\texttt{RES-0153}}] Two matrix paths with the same endpoint configuration can do different work on fermions. Matter travelling near the same incoming ray can therefore retain a distinction missed at the endpoints.
\item[\textnormal{\texttt{RES-0154}}] A timelike material matrix block supplies an oriented three-plane and absorbs work. Excited levels and an individual signed charge crossing can occur in one dynamical solution.
\item[\textnormal{\texttt{RES-0155}}] A Yukawa interaction between two blocks transports a compensating charge. Finite energy constrains its leakage and the probability that the current record becomes neutral.
\item[\textnormal{\texttt{RES-0156}}] Half-filled material states can receive charge while preserving an unknown helicity pair conditional on the recorded group. Recording a microscopic spin direction would change that conclusion.
\item[\textnormal{\texttt{RES-0157}}] Relative motion can pay the energy needed for a receipt. Finite material spin limits both angular resolution and the orbital information that survives a high-curvature passage.
\item[\textnormal{\texttt{RES-0158}}] The falling receiver has physical polarizations and can exchange charge without prescribing an arriving quantum by hand in the example. Local orbits can distinguish two candidate interiors.
\item[\textnormal{\texttt{RES-0159}}] In one ten-matrix solution two deformations create pairs from the vacuum and drive a charge response. A fixed finite-energy wave packet limits how much material information an increasingly sharp tide carries.
\item[\textnormal{\texttt{RES-0160}}] A moving matrix sphere has a geometrically defined area; supporting a prescribed central motion may require extra force. Compensated transfer and recoil can occur before its area vanishes.
\item[\textnormal{\texttt{RES-0161}}] Specific internal matrix data can match Schwarzschild tidal components on an interior slice. A subsequent bounce and quantum lower bound for area depend on the incoming quantum state.
\item[\textnormal{\texttt{RES-0162}}] A curved vacuum Einstein wave coexists with matrix bounces requiring negative radial null convergence. Different energies can discriminate continuations. Local matrix dynamics and global smooth structure are separate issues.
\item[\textnormal{\texttt{RES-0163}}] For a rotating example angular momentum bounds both spatial area and volume. Curved private patches can share an open region; global continuation, a physical longitudinal momentum and a layer label have distinct tests.
\item[\textnormal{\texttt{RES-0164}}] The Weyl tensor at a rotating bounce, oriented crossings and material compensation can be compared in one preparation. Finite-rank matter and a D2 magnetic probe need not resolve the same distinctions.
\item[\textnormal{\texttt{RES-0165}}] An ellipsoidal branch gives regular matter propagation on its symmetry axis. Absorption of a graviton with compensating recoil, measurable shear and finite resolution on the M circle supply independent tests.
\item[\textnormal{\texttt{RES-0166}}] The sign of period-integrated null convergence, a fifth-angular-mode threshold and an extra optical arm distinguish branches. Signed outgoing matter counts can resolve M-momentum transfer that total counts miss.
\item[\textnormal{\texttt{RES-0167}}] A freely falling elastic sphere can resolve selected components of the tide. Finite bandwidth, empty gravitational contacts and a short material response deadline constrain irreversible directional records.
\end{description}

\subsection{Vacua, astronomical tests, the internal radius and forces on matter}
\begin{description}\small
\setlength{\itemsep}{.25em}\setlength{\parskip}{0pt}
\item[\textnormal{\texttt{RES-0168}}] In the specified orientifold catalogue the symmetry quotient gives 115 flux orbits, 49 linearly fixing ten complex variables. Eleven astronomical tests distinguish selected effective branches up to cosmological distances.
\item[\textnormal{\texttt{RES-0169}}] Among the specified pairs of flux configurations, most share only a scale deformation. A six-polarization local emission can preserve a four-component quantum code jointly in radiation and residual material.
\item[\textnormal{\texttt{RES-0170}}] Several layers can agree on flux marginals and disagree on permitted shapes. A falling magnetized material develops a finite endpoint current; its separate endpoint counts do not certify the required routing charge.
\item[\textnormal{\texttt{RES-0171}}] Primitive Euclidean branes can exclude an internal square point even in flux sectors with fixed complex shape. A conditional local minimum in four complex fields leaves open the full radius and spectrum problem.
\item[\textnormal{\texttt{RES-0172}}] Rotating matrix matter can stabilize some deformations and develop an axial or finite-wavelength instability in others. A stationary radius in a truncation does not establish a stable four-dimensional vacuum.
\item[\textnormal{\texttt{RES-0173}}] A tetrahedral projection gives a fixed-charge minimum for one open-string species. Additional string species or a coupling to clocks can destabilize it or make its size observable at long distances.
\item[\textnormal{\texttt{RES-0174}}] A clock with fixed proper volume and charge can retain ordinary gravitational redshift. Finite exchange, surface stress and one-loop tensor propagation separately constrain a proposed confining mechanism.
\item[\textnormal{\texttt{RES-0175}}] A localized charged matrix object can have a restoring force for its outer radius. Its external size and the scale of the internal space are distinct dynamical degrees of freedom.
\item[\textnormal{\texttt{RES-0176}}] Allowing the four-dimensional metric to vary leaves a common-scale direction in a specified reduction. Four rotating links can reproduce a matter wave while retaining a negative physical mode.
\item[\textnormal{\texttt{RES-0177}}] In a three-link magnetic projection, a positive two-spin gravitational response requires a negative bare mass-squared in the tested branch. Rotational support trades range against force and energy density.
\item[\textnormal{\texttt{RES-0178}}] Certain quantized type-II sources can supply restoring forces for 38 modes in a restricted model. Finite-wavelength instabilities and the electromagnetic dependence of atomic masses constrain this route.
\item[\textnormal{\texttt{RES-0179}}] Three types of magnetic charge can fix a 70-scalar truncation locally in one configuration; a second configuration leaves eight exact directions. Chiral light matter imposes a separate test.
\item[\textnormal{\texttt{RES-0180}}] A charge balance can fix three radii yet leave a long-range composition-dependent atomic force. In the light unscreened electromagnetic example the predicted Ti/Pt signal is far above the measured bound.
\item[\textnormal{\texttt{RES-0181}}] On the same static six-dimensional background ordinary CAT, IKKT and static BFSS share a one-loop cancellation in the tested sector. Freely populated positive-energy bulk modes cannot oppose the common-scale force there.
\end{description}

\subsection{A common source, physical output domains and a test of reversal}
\begin{description}\small
\setlength{\itemsep}{.25em}\setlength{\parskip}{0pt}
\item[\textnormal{\texttt{RES-0182}}] One prepared polarized source gives positively confirmed A and an exactly closed B sector with probability equal to the source weight times the cube of the bright absorption probability. A recoil-conserving boundary encodes one unknown qubit in orthogonal radiation or core carriers. Inward histories erase the qubit for B.
\item[\textnormal{\texttt{RES-0183}}] The real-time scalar inherits positive clock energy. One network has the same inclusive channel for either reader, while a physically selected port can change the inward dimension. Both readout signs and a uniform transfer bound are needed to test the return probability.
\end{description}
